\documentclass[11pt,a4paper]{article}

\usepackage[margin=1.8cm]{geometry}
\usepackage{fontspec}
\usepackage{enumitem}
\usepackage{xurl}
\usepackage[hidelinks]{hyperref}
\usepackage{xcolor}
\usepackage{titlesec}

\setmainfont{Times New Roman}
\definecolor{heading}{HTML}{003366}
\definecolor{subtle}{HTML}{555555}

\pagestyle{empty}
\setlength{\parindent}{0pt}
\setlength{\parskip}{5pt}
\sloppy
\setlist[itemize]{leftmargin=1.35em, topsep=2pt, itemsep=2pt}

\titleformat{\section}{\large\bfseries\color{heading}}{}{0pt}{}[\titlerule]
\titleformat{\subsection}{\normalsize\bfseries\color{heading}}{}{0pt}{}

\begin{document}

{\LARGE \textbf{Research Statement}}\\[-1pt]
\textbf{Wen Quan}\\
\textcolor{subtle}{Proposed Research Direction: AI-Assisted Assessment and Design Optimization for Age-Friendly Built Environments}\\
\href{mailto:wenquan@student.usm.my}{wenquan@student.usm.my} \,|\, 
\href{https://wenquan0816.github.io/schoolar-site/}{\nolinkurl{wenquan0816.github.io/schoolar-site}}

\section*{Research Vision}
My research develops artificial intelligence methods for evaluating, explaining, and improving age-friendly built environments. I am especially interested in how AI can translate complex evidence from buildings, communities, policies, user needs, and spatial data into design decisions that support older adults' safety, accessibility, health, autonomy, and dignity.

The central question guiding my next stage of work is: \textit{How can AI-assisted environmental assessment become a reliable design and decision-support system for age-friendly homes, communities, care facilities, and public environments?} This question sits between interior design, architecture, built environment research, public health, gerontology, and computational methods. It also aligns with Malaysia's need for inclusive housing, healthy ageing, universal design, and digital transformation in the built environment.

My publication record has already established this direction. I have developed AI-enabled assessment frameworks for age-friendly living environments, public health emergency resilience, rural school-to-elderly-care conversion, and building energy flexibility. These studies use methods such as natural language processing, contrastive learning, hyperbolic embedding, LSTM-based model predictive control, SciBERT-assisted indicator extraction, expert validation, and multi-dimensional assessment. My future work will consolidate these into a coherent research program on \textbf{AI+Age-Friendly Environment Assessment}.

\section*{Research Agenda}
\subsection*{1. AI-Based Indicator Systems for Age-Friendly Environmental Assessment}
The first strand of my research will build robust indicator systems for assessing age-friendly environments across scales: room, home, residential community, care facility, and public space. Traditional assessment frameworks often depend on static checklists and expert judgement. My approach combines AI-assisted evidence extraction with domain validation:
\begin{itemize}
  \item using NLP and large language models to extract candidate indicators from standards, policies, interview transcripts, design guidelines, and academic literature;
  \item applying Delphi, fuzzy Delphi, AHP, entropy weighting, or explainable machine learning to validate and weight indicators;
  \item building assessment tools that classify environmental problems into actionable design priorities.
\end{itemize}

This strand will support practical outputs such as age-friendly renovation checklists, care-facility conversion evaluation tools, and design decision matrices for housing and community renewal. It can also generate publishable work in built environment, gerontology, public health, and design research outlets.

\subsection*{2. Multimodal AI for Spatial Accessibility, Safety, and Comfort}
The second strand will develop multimodal assessment models that combine visual, spatial, textual, and sensor-based data. Many age-friendly design problems are spatial and physical: step heights, doorway widths, circulation paths, handrail locations, lighting quality, thermal comfort, furniture layout, fall-risk zones, and access to community facilities. These can be assessed more consistently with AI if the data pipeline is carefully designed.

My planned methods include:
\begin{itemize}
  \item computer vision and 3D point-cloud analysis for detecting accessibility barriers and fall-risk features;
  \item BIM- and digital-twin-linked assessment for comparing existing spatial conditions with age-friendly design requirements;
  \item thermal and environmental comfort models for elderly day-care and residential settings;
  \item explainable AI methods, such as SHAP and attention visualization, to make model results understandable to designers, care providers, and policy stakeholders.
\end{itemize}

The aim is not to replace professional judgement, but to provide designers and decision-makers with faster, evidence-based screening tools. For a Malaysian university context, this strand can connect with architecture, interior architecture, building technology, data science, and public health collaborators.

\subsection*{3. AI-Supported Design Optimization and Renovation Decision-Making}
The third strand will move from assessment to design optimization. Once environmental risks and deficiencies are identified, the next question is which design interventions should be prioritized under budget, technical, cultural, and operational constraints. I plan to develop AI-supported workflows for:
\begin{itemize}
  \item prioritizing renovation actions for older residential communities and long-term care facilities;
  \item generating alternative age-friendly interior design and layout solutions;
  \item evaluating trade-offs among safety, accessibility, cost, comfort, and user preference;
  \item linking environmental assessment with policy and facility-management decisions.
\end{itemize}

This direction builds on my PI-led projects on aging-in-place renovation, universal design for new communities, hospital outpatient interior environments, community-based senior living, and adaptive renovation of long-term care facilities under COVID-19 protocols. It provides a clear pathway for grant proposals and industry/community collaboration.

\section*{Three-Year Research Plan}
\textbf{Year 1: Consolidation and local collaboration.} I will refine my existing AI-assisted assessment frameworks into a Malaysian-applicable research program, build partnerships with colleagues in interior architecture, architecture, building technology, planning, gerontology, and public health, and prepare pilot studies on age-friendly housing/community environments. Outputs will include a scoping review, a validated indicator framework, and one or two grant applications.

\textbf{Year 2: Data collection and model development.} I will collect multimodal data from selected residential, community, or care environments, including policy/design documents, environmental observations, spatial measurements, photographs or 3D scans where ethically appropriate, and stakeholder interviews. I will develop AI models for indicator extraction, environmental classification, and spatial risk detection, with careful attention to data governance, consent, privacy, and explainability.

\textbf{Year 3: Design translation and applied impact.} I will translate the models into practical assessment and design-support tools, such as an age-friendly environment assessment dashboard, renovation prioritization matrix, or BIM-linked evaluation workflow. I will target publications in WoS/Scopus-indexed journals and pursue collaboration with local authorities, care providers, housing stakeholders, or design practices.

\section*{Fit with Malaysian University Roles}
This research program is well suited to lecturer roles in Interior Design, Interior Architecture, Architecture, Built Environment, Environmental Design, Design Technology, and Age-Friendly Design. It can support studio teaching, research-led design projects, postgraduate supervision, external grants, and interdisciplinary collaboration. It also offers a distinctive contribution: connecting AI methods with socially responsive design for ageing societies.

\end{document}
